\begin{ouverture}[statistiques]{20}{Lire et croiser des données}{Partie 5 · Modéliser et organiser}
{Tableaux · Barres · Courbes · Diagrammes circulaires}
\objectif{Lire tableaux, diagrammes en barres, courbes et diagrammes circulaires.}
\objectif{Construire un tableau ou un diagramme à partir de données.}
\objectif{Croiser plusieurs documents pour répondre à une question.}
\prerequis{Lectures de graphiques du CM1.}
\end{ouverture}
\pagestyle{otMain}\essentiel\fontsize{11}{15}\selectfont
\phantomsection\label{cours-20-1}
\essSection{1}{Des tableaux aux diagrammes}
\begin{definition}Un tableau à double entrée se lit au croisement d’une ligne et d’une colonne. Un diagramme en barres porte un effectif par barre, lu sur l’axe gradué. Les totaux contrôlent l’ensemble : la somme des effectifs retrouve le nombre de réponses.\end{definition}
\begin{exemple}Les données réelles, souvent liées au climat ou à la biodiversité, s’annoncent avec leurs unités : millimètres de pluie, degrés Celsius, nombre d’espèces. Les lire, c’est d’abord lire le titre et les axes.\end{exemple}
\begin{savoirfaire}[Le pas des graduations]Avant toute lecture : repérer le pas des graduations. Un carreau peut valoir 2, 5, 10 ou 100 : le même dessin raconte alors d’autres nombres.\end{savoirfaire}
\begin{attention}Entre deux points mesurés, une courbe relie sans garantir : les valeurs intermédiaires sont des estimations, pas des mesures.\end{attention}
\clearpage
\phantomsection\label{cours-20-2}
\essSection{2}{Le diagramme circulaire}
\begin{definition}Un diagramme circulaire partage un disque proportionnellement aux effectifs : la moitié du disque représente la moitié du total, un quart en représente le quart. On le lit en comparant les secteurs au disque entier.\end{definition}
\label{fig-cours-20-2-cours-circulaire}\figureGeo{cours-circulaire}

\begin{savoirfaire}[Du secteur à l’effectif]Lire un circulaire : identifier le total, puis traduire chaque secteur en fraction simple du disque : moitié, quart, tiers. Un secteur d’un quart sur 60 réponses représente $60:4=15$ réponses.\end{savoirfaire}
\begin{attention}Le circulaire montre des parts du tout, pas une évolution : pour suivre une grandeur dans le temps, la courbe reste l’outil adapté. Choisir le graphique selon la question posée.\end{attention}

\clearpage\rubrique{Découverte}
\begin{blocLivre}{ex-20-1}\begin{exercice}[1]{Le tableau de l’enquête}
L’enquête « sport préféré » donne :\begin{center}\renewcommand{\arraystretch}{1.35}\begin{tabular}{c|cccc}\toprule
Sport & foot & danse & judo & natation\\\midrule
Effectif & 12 & 9 & 5 & 4\\
\bottomrule\end{tabular}\end{center}Combien d’élèves ont répondu ? Quel sport arrive en tête ?
\end{exercice}\end{blocLivre}
\begin{blocLivre}{ex-20-2}\begin{exercice}[1]{Double entrée}
Le tableau croise deux classes et deux sorties.\begin{center}\renewcommand{\arraystretch}{1.35}\begin{tabular}{c|cc}\toprule
 & Musée & Forêt\\\midrule
CM2 A & 15 & 11\\
CM2 B & 9 & 17\\
\bottomrule\end{tabular}\end{center}Combien d’élèves de CM2 B préfèrent la forêt ? Combien d’élèves préfèrent le musée ?
\end{exercice}\end{blocLivre}
\begin{blocLivre}{ex-20-3}\begin{exercice}[1]{Lire les barres}
Le diagramme donne les livres empruntés chaque jour ; une graduation vaut 5. Lire les valeurs de lundi et de jeudi, et dire quel jour bat le record.
\label{fig-ex-20-3}\figureGeo{don-3}
\end{exercice}\end{blocLivre}
\begin{blocLivre}{ex-20-4}\begin{exercice}[1]{La courbe de l’année}
La courbe donne la température moyenne de chaque mois. Quel est le mois le plus froid ? Quel est l’écart entre juillet et janvier ?
\label{fig-ex-20-4}\figureGeo{don-4}
\end{exercice}\end{blocLivre}
\begin{blocLivre}{ex-20-5}\begin{exercice}[1]{Le circulaire de la cantine}
Le diagramme circulaire montre les choix de 60 élèves : la moitié a choisi le plat A, un quart le plat B, le reste le plat C. Donner les trois effectifs.
\label{fig-ex-20-5}\figureGeo{don-5}
\end{exercice}\end{blocLivre}

\rubrique{Entraînement}
\begin{blocLivre}{ex-20-6}\begin{exercice}[2]{Construire le diagramme}
Construire le diagramme en barres des sports de l’exercice 1, avec une graduation de 2.
\label{fig-ex-20-6}\figureGeo{don-6}
\end{exercice}\end{blocLivre}
\begin{blocLivre}{ex-20-7}\begin{exercice}[2]{Compléter le tableau}
Compléter les cases manquantes :\begin{center}\renewcommand{\arraystretch}{1.35}\begin{tabular}{c|ccc}\toprule
 & Lundi & Mardi & Total\\\midrule
Entrées & 145 & ? & 313\\
Goûters & ? & 96 & 184\\
\bottomrule\end{tabular}\end{center}
\end{exercice}\end{blocLivre}
\begin{blocLivre}{ex-20-8}\begin{exercice}[2]{Les points du repère}
Chaque point donne la taille d’un plant de tomate selon la semaine. Lire la taille en semaine 2 et en semaine 5, puis dire entre quelles semaines la croissance a été la plus forte.
\label{fig-ex-20-8}\figureGeo{don-8}
\end{exercice}\end{blocLivre}
\begin{blocLivre}{ex-20-9}\begin{exercice}[2]{Vrai ou faux sur le circulaire}
Sur un circulaire, le secteur « vélo » couvre un quart du disque. Peut-on dire que 25 élèves viennent à vélo ? Justifier.
\end{exercice}\end{blocLivre}
\begin{blocLivre}{ex-20-10}\begin{exercice}[2]{Deux documents à croiser}
Le tableau donne les ventes du matin : 48 pains ; la courbe indique 76 pains vendus sur la journée. Combien de pains ont été vendus l’après-midi ?
\label{fig-ex-20-10}\figureGeo{don-10}
\end{exercice}\end{blocLivre}

\rubrique{Problèmes}
\begin{blocLivre}{ex-20-11}\begin{exercice}[3]{La barre incohérente}
Le tableau annonce : CP 22, CE1 26, CE2 24. Une barre du diagramme contredit le tableau : laquelle ? Corriger.
\label{fig-ex-20-11}\figureGeo{don-11}
\end{exercice}\end{blocLivre}
\begin{blocLivre}{ex-20-12}\begin{exercice}[3]{L’enquête à construire}
Réponses brutes à « Combien de livres as-tu lus ce mois-ci ? » : 0, 1, 2, 1, 0, 3, 1, 2, 1, 0, 2, 1.\begin{questions}
\item Construire le tableau des effectifs.
\item Combien d’élèves ont lu au moins 2 livres ?
\end{questions}
\end{exercice}\end{blocLivre}
\begin{blocLivre}{ex-20-13}\begin{exercice}[3]{Le climat en deux graphiques}
Pour une même ville, la courbe donne les températures mensuelles et le diagramme en barres les précipitations. Quel mois cumule chaleur maximale et pluie minimale ?
\label{fig-ex-20-13}\figureGeo{don-13}
\end{exercice}\end{blocLivre}
\begin{blocLivre}{ex-20-14}\begin{exercice}[3]{Choisir le bon graphique}
Pour chaque question, choisir l’outil adapté - tableau, courbe ou circulaire : suivre la taille d’une plante semaine après semaine ; montrer la part de chaque sport dans la classe ; retrouver une valeur exacte parmi cinquante relevés.
\end{exercice}\end{blocLivre}
\begin{blocLivre}{ex-20-15}\begin{exercice}[3]{Le piège de la lecture rapide}
Un diagramme en barres sur les ventes commence ses graduations à 50, et la barre B semble deux fois plus haute que la barre A. A : 60 ventes, B : 70 ventes. Expliquer l’illusion.
\end{exercice}\end{blocLivre}

\clearpage
